\section{Matriz de trazabilidad de escenarios}

\begin{longtable}{p{1.8cm}p{3.2cm}p{4.4cm}p{2.7cm}}
\caption{Trazabilidad requisito--actividad--evidencia}\label{tab:trace}\\
\toprule
Escenario & Actividad & Evidencia principal & Condición de origen \\
\midrule
\endfirsthead
\toprule
Escenario & Actividad & Evidencia principal & Condición de origen \\
\midrule
\endhead
1 & Análisis del caso & \texttt{docs/}\allowbreak\texttt{analisis\_inicial.md} & Literal \\
2 & Clasificación ISO 25022 & CSV clasificado y resumen JSON & Literal \\
3 & Niveles de calidad y SQuaRE & Mapa, SonarQube, JMeter y Figura \ref{fig:square} & Literal \\
4 & Atributos UTC & Fichas \ref{tab:utc-hce}--\ref{tab:utc-facturacion} y Tabla \ref{tab:utc} & Guiado / aplicado \\
5 & Mapeo y prioridad & Matriz \ref{tab:priority} & Guiado / aplicado \\
6 & Diseño de métricas & Fichas, fórmulas y Tabla \ref{tab:catalog} & Guiado / aplicado \\
7 & Obtención de datos & Tablas \ref{tab:event-types} y \ref{tab:data-validation}; manifiesto & Guiado / aplicado \\
8 & Automatización & Figura \ref{fig:pipeline-phase8}, 19 pruebas ejecutadas y flujo CI & Guiado / aplicado \\
9 & KPI y tablero & Tabla \ref{tab:results} y React gerencial & Reconstruido \\
10 & Interpretación & Hipótesis de causa y prioridades & Reconstruido \\
11 & Comunicación ejecutiva & Síntesis de decisiones & Reconstruido \\
12 & Mejora continua & Plan PDCA, Tabla \ref{tab:pdca} & Reconstruido \\
Reto & Telemedicina 2.0 & Trazabilidad y casos de contexto & Reconstruido \\
\bottomrule
\end{longtable}

\section{Ubicación de artefactos}

\begin{table}[H]
\centering
\caption{Estructura de evidencia reproducible}
\begin{tabularx}{\textwidth}{p{4.0cm}Y}
\toprule
Ruta & Responsabilidad \\
\midrule
\texttt{docs/}\allowbreak\texttt{analisis\_inicial.md} & Matriz de análisis solicitada en el escenario 1. \\
\texttt{data/}\allowbreak\texttt{incidentes\_2025.csv} & Dataset solicitado por la actividad, ampliado a 3.000 registros. \\
\texttt{data/clasificacion} & Tabla 2.2 completada para los 3.000 incidentes, incluidos 1001--1006. \\
\texttt{apps/data/original} & Copia preservada del dataset de 3.000 incidentes. \\
\texttt{apps/data/simulated} & Eventos, encuestas, manifiestos y ejecuciones regenerables. \\
\texttt{apps/data/processed} & Clasificación y métricas derivadas. \\
\texttt{apps/analytics} & Clasificación, simulación, medición y exportación en Python. \\
\texttt{apps/mobile-patient} & Portal Flutter del paciente. \\
\texttt{apps/web-his} & HIS operativo React y módulo gerencial de calidad. \\
\texttt{apps/backend} & Gateway, dominios clínicos y API de medición. \\
\texttt{reportes/evidencias/}\allowbreak\texttt{resultados} & Resúmenes seleccionados para entrega. \\
\texttt{reportes/evidencias/}\allowbreak\texttt{calidad} & Métricas SonarQube y resumen de la prueba JMeter. \\
\texttt{reportes/evidencias/}\allowbreak\texttt{pruebas} & Acta Markdown y resumen JSON de las 19 pruebas ejecutadas el 13 de julio de 2026. \\
\texttt{reportes/entrega} & PDF académico final. \\
\bottomrule
\end{tabularx}
\end{table}

\section{Controles de reproducción}

\begin{enumerate}
  \item Verificar el SHA-256 registrado de \texttt{apps/}\allowbreak\texttt{data/}\allowbreak\texttt{original/}\allowbreak\texttt{incidentes\_2025.csv} sin modificar ese archivo.
  \item Instalar las dependencias de Python y ejecutar el pipeline con semilla 25022.
  \item Confirmar 3.000 clasificaciones, 6.000 eventos, 2.000 encuestas y 10 métricas.
  \item Repetir la simulación y comparar manifiestos y resultados.
  \item Levantar Docker Compose desde Ubuntu WSL y comprobar los estados de salud.
  \item Ejecutar SonarQube sobre facturación y exportar sus métricas mediante \texttt{run-sonar-wsl.sh}.
  \item Ejecutar JMeter con 500 usuarios y generar el resumen mediante \texttt{run-jmeter-wsl.sh}.
  \item Probar React y Flutter contra el mismo Gateway local.
  \item Compilar este documento con XeLaTeX mediante \texttt{latexmk}.
  \item Renderizar todas las páginas y revisar portada, tablas, figuras, referencias y desbordamientos.
\end{enumerate}

\section{Acta de ejecución de pruebas}

El 13 de julio de 2026 se ejecutaron de forma efectiva seis suites con 19 casos: analítica Python (5), API de medición (1), dominio HCE Java (1), portal React (8), integración Docker (3) y portal Flutter (1). Los 19 casos aprobaron. También finalizaron correctamente el pipeline de diez métricas, el análisis estático Flutter, las construcciones Vite y Flutter Web, el levantamiento saludable de la pila y una carga JMeter de 500 solicitudes sin errores.

La evidencia detallada conserva comandos, versiones, tiempos, incidencias y resolución. El acta y su resumen estructurado están en \texttt{reportes/evidencias/pruebas}; la evidencia dinámica de rendimiento, incluido el archivo JTL, está en \texttt{reportes/evidencias/calidad/jmeter}.

\section{Criterio de procedencia}

Para evitar ambigüedad, todo resultado debe identificarse con una de estas etiquetas:

\begin{description}[style=nextline,leftmargin=2.5cm,labelwidth=2.2cm]
  \item[Original:] contenido recibido en la guía del taller o en el dataset suministrado.
  \item[Simulado:] observación generada artificialmente con semilla y reglas declaradas.
  \item[Derivado:] clasificación, agregado, fórmula, gráfica o interpretación calculada a partir de una fuente.
  \item[Propuesto:] métrica, meta, hipótesis o acción definida por el equipo para completar el taller.
\end{description}
